\section*{Chapter \ref{ch:g12:limcont} --- Limits and Continuity}

\begin{solution}{exo:g12:limcont:1}
Dividing by $x^3$:
$\dfrac{2x^3 - x + 1}{5x^3 + x^2} = \dfrac{2 - 1/x^2 + 1/x^3}{5 + 1/x}
\to \dfrac{2}{5}$.

$x^2 + 3x - 1 = x^2(1 + 3/x - 1/x^2) \to +\infty$ as $x \to -\infty$, since
$x^2 \to +\infty$ and the bracket tends to $1$.

As $x \to 2^+$, the numerator tends to $3 > 0$ and the denominator to $0^+$,
so $\dfrac{x+1}{x-2} \to +\infty$.

For $x > 0$, $-\dfrac1x \leq \dfrac{\cos x}{x} \leq \dfrac1x$, so the limit
is $0$ by the squeeze theorem.
\end{solution}

\begin{solution}{exo:g12:limcont:2}
\emph{1.} $2x^2 - 3x + 1 = (x-1)(2x-1)$ (the \omterm{def:g11:quad:discriminant}{root} $x=1$ is apparent from
$2 - 3 + 1 = 0$). For $x \neq 1$, cancel $x - 1$: $f(x) = 2x - 1$.

\emph{2.} Hence $\lim_{x\to1} f(x) = 1$. Setting $f(1) = 1$ extends $f$ to
the polynomial \omterm{def:g11:func:function}{function} $x \mapsto 2x-1$ on $\R$, which is \omterm{def:g12:limcont:continuity}{continuous} at $1$.
\end{solution}

\begin{solution}{exo:g12:limcont:3}
Write $f(x) = \dfrac{3x - 1}{x + 2} = \dfrac{3(x+2) - 7}{x+2}
= 3 - \dfrac{7}{x+2}$. As $x \to \pm\infty$, $f(x) \to 3$: the line $y = 3$
is a \omterm{def:g12:limcont:asymptote}{horizontal asymptote} at both infinities. As $x \to (-2)^\pm$,
$f(x) \to \mp\infty$: the line $x = -2$ is a \omterm{def:g12:limcont:asymptote}{vertical asymptote}. The curve is
a \omterm{def:g11:func:reference}{hyperbola}, below $y=3$ for $x > -2$ and above it for $x < -2$.
\end{solution}

\begin{solution}{exo:g12:limcont:4}
Multiply by the conjugate:
\[
\sqrt{x^2 + x} - x = \frac{x^2 + x - x^2}{\sqrt{x^2+x} + x}
= \frac{x}{x\left(\sqrt{1 + 1/x} + 1\right)}
= \frac{1}{\sqrt{1 + 1/x} + 1}
\xrightarrow[x\to+\infty]{} \frac{1}{2}.
\]
Similarly,
\[
\frac{\sqrt{1+x} - 1}{x} = \frac{(1+x) - 1}{x\left(\sqrt{1+x} + 1\right)}
= \frac{1}{\sqrt{1+x} + 1} \xrightarrow[x\to0]{} \frac{1}{2}.
\]
\end{solution}

\begin{solution}{exo:g12:limcont:5}
\emph{1.} Since $\sin x \geq -1$, $f(x) \geq x - 1 \to +\infty$, and we
conclude by comparison. Yet $f'(x) = 1 + \cos x$ vanishes at every odd
multiple of $\pi$, and on every \omterm{def:g10:numbers:interval}{interval} $\intco{A}{+\infty}$ the \omterm{def:g11:func:function}{function}
alternates between \omterm{def:g10:numbers:interval}{intervals} of strict increase; it is \omterm{def:g12:seq:monotonic}{increasing} but not
strictly, and more importantly its growth is not needed: comparison suffices.

\emph{2.} For $x > 1$, divide by $x$:
\[
g(x) = \frac{1 + \frac{\sin x}{x}}{1 - \frac{\sin x}{x}},
\]
and $\frac{\sin x}{x} \to 0$ by the squeeze theorem
($\abs{\sin x} \leq 1$). Hence $g(x) \to \frac{1+0}{1-0} = 1$.
\end{solution}

\begin{solution}{exo:g12:limcont:6}
$f(x) = x^3 - 3x + 1$ is \omterm{def:g12:limcont:continuity}{continuous}, with $f'(x) = 3x^2 - 3 = 3(x-1)(x+1)$:
$f$ increases on $\intoc{-\infty}{-1}$, decreases on $\intcc{-1}{1}$,
increases on $\intco{1}{+\infty}$, with local \omterm{def:g10:functions:extrema}{maximum} $f(-1) = 3 > 0$ and
local \omterm{def:g10:functions:extrema}{minimum} $f(1) = -1 < 0$, $\lim_{-\infty} f = -\infty$,
$\lim_{+\infty} f = +\infty$.

On each of the three monotonicity \omterm{def:g10:numbers:interval}{intervals}, $f$ is \omterm{def:g12:limcont:continuity}{continuous}, strictly
\omterm{def:g12:seq:monotonic}{monotonic}, and $0$ lies strictly between the endpoint values/limits, so the
bijection theorem gives exactly one zero per \omterm{def:g10:numbers:interval}{interval}: three zeros in total.
Sign checks: $f(-2) = -1 < 0 < f(-1) = 3$, so one zero in
$\intoo{-2}{-1}$; $f(0) = 1 > 0 > f(1) = -1$, one zero in $\intoo{0}{1}$;
$f(1) = -1 < 0 < f(2) = 3$, one zero in $\intoo{1}{2}$.
\end{solution}

\begin{solution}{exo:g12:limcont:7}
Let $g(x) = f(x) - x$, \omterm{def:g12:limcont:continuity}{continuous} on $\intcc{0}{1}$. Since
$f(0) \in \intcc{0}{1}$, $g(0) = f(0) \geq 0$; since $f(1) \in \intcc{0}{1}$,
$g(1) = f(1) - 1 \leq 0$. By the intermediate value theorem, $g$ vanishes at
some $c \in \intcc{0}{1}$, i.e.\ $f(c) = c$.
\end{solution}

\begin{solution}{exo:g12:limcont:8}
Let $u(t)$ and $d(t)$ be the positions (measured as distance from the bottom
of the trail) at time $t \in \intcc{8}{20}$ on the way up and on the way
down. Both are \omterm{def:g12:limcont:continuity}{continuous}, $u(8) = 0$, $u(20) = L$ (the length of the trail),
$d(8) = L$, $d(20) = 0$. The \omterm{def:g11:func:function}{function} $g = u - d$ is \omterm{def:g12:limcont:continuity}{continuous} with
$g(8) = -L \leq 0$ and $g(20) = L \geq 0$, so $g(c) = 0$ for some time $c$ by
the intermediate value theorem: at time $c$, both days' positions coincide.
\end{solution}

\begin{solution}{exo:g12:limcont:9}
\emph{1.} $f'(x) = 5x^4 + 3x^2 + 1 > 0$, so $f$ is strictly \omterm{def:g12:seq:monotonic}{increasing} on
$\R$; with $\lim_{\pm\infty} f = \pm\infty$ and \omterm{def:g12:limcont:continuity}{continuity}, the bijection
theorem gives a unique real zero $\alpha$. Since $f(0) = -1 < 0$ and
$f(1) = 2 > 0$, $\alpha \in \intoo{0}{1}$.

\emph{2.} Maintain an \omterm{def:g10:numbers:interval}{interval} $\intcc{a_n}{b_n}$ with
$f(a_n) < 0 < f(b_n)$, starting from $\intcc{0}{1}$; at each step evaluate
$f$ at the \omterm{prop:g10:coordgeom:midpoint}{midpoint} $m$ and keep the \omterm{def:g10:numbers:interval}{half-interval} on which the sign change
persists. After $n$ steps the \omterm{def:g10:numbers:interval}{interval} has length $2^{-n}$, so the \omterm{prop:g10:coordgeom:midpoint}{midpoint}
approximates $\alpha$ to within $2^{-(n+1)}$. We need
$2^{-(n+1)} \leq 10^{-6}$, i.e.\ $n + 1 \geq \frac{6}{\log_{10} 2} \approx 19.93$:
$n = 20$ iterations suffice.
\end{solution}

\begin{solution}{pb:g12:limcont:1}
\textbf{1.} Dominant terms: $\frac{3x^2}{x^2} \to 3$;
$\frac{2x}{x^2} = \frac2x \to 0$;
$x^3 - x^2 = x^3\left(1 - \frac1x\right) \to +\infty$.

\textbf{2.} $\lim_{x \to \pm\infty} f = 3$: \omterm{def:g12:limcont:asymptote}{horizontal
asymptote} $y = 3$; $\lim_{x \to (-2)^\pm} f = \mp\infty$:
\omterm{def:g12:limcont:asymptote}{vertical asymptote} $x = -2$.

\textbf{3.} $\frac{(x-3)(x+3)}{x - 3} = x + 3 \to 6$. And
$\frac{\sqrt{x+1} - 1}{x} = \frac{x}{x(\sqrt{x+1} + 1)} =
\frac{1}{\sqrt{x+1} + 1} \to \frac12$.

\textbf{4.} $\abs{x \sin\frac1x} \leq \abs x \to 0$, so the
limit is $0$. And
$\abs{\frac{\sin x}{x}} \leq \frac1x \to 0$: limit $0$.

\textbf{5.} $\frac{\sqrt{4x^2 + x}}{x} =
\sqrt{4 + \frac1x} \to \sqrt4 = 2$ (\omterm{def:g12:limcont:continuity}{continuity} of the square
root, composition).

\textbf{6.} $f(x) = x^3 + x$ is \omterm{def:g12:limcont:continuity}{continuous} with
$f(1) = 2 < 5 < 10 = f(2)$: by the IVT, $f(c) = 5$ for some
$c \in \intoo{1}{2}$.

\textbf{7.} $f'(x) = 3x^2 + 1 > 0$: strictly \omterm{def:g12:seq:monotonic}{increasing} on
$\R$, \omterm{def:g12:limcont:continuity}{continuous}, with limits $\mp\infty$: by the bijection
theorem, every value --- in particular $5$ --- is attained
exactly once.

\textbf{8.} $p(x) = x^3\left(1 + \frac ax + \frac{b}{x^2} +
\frac{c}{x^3}\right)$: as $x \to +\infty$, $p \to +\infty$; as
$x \to -\infty$, $p \to -\infty$. So on some huge \omterm{def:g10:numbers:interval}{interval}
$\intcc{-M}{M}$, $p(-M) < 0 < p(M)$: the IVT delivers a \omterm{def:g11:quad:discriminant}{root}.
The same dominant-term argument works verbatim for any odd
degree --- odd powers keep opposite signs at the two
infinities.

\textbf{9.} Halving $n$ times leaves an \omterm{def:g10:numbers:interval}{interval} of length
$2^{-n}$; $2^{-20} \approx 9.5 \times 10^{-7} < 10^{-6}$:
twenty steps. Newton doubles correct digits each step: from
one correct digit, about $3$--$4$ steps reach six --- the
tortoise and the greyhound of root-finding.

\textbf{10.} $g$ is not \omterm{def:g12:limcont:continuity}{continuous} on $\intcc{-1}{1}$: it
explodes at $0$, which the \omterm{def:g10:numbers:interval}{interval} contains. The IVT's
conclusion evaporates with its hypothesis --- always check
\omterm{def:g12:limcont:continuity}{continuity} \emph{on the whole \omterm{def:g10:numbers:interval}{interval}}, not just at the
endpoints.

\textbf{11.} $g(x) = f(x) - x$ is \omterm{def:g12:limcont:continuity}{continuous};
$g(0) = f(0) \geq 0$ (values lie in $\intcc{0}{1}$) and
$g(1) = f(1) - 1 \leq 0$. By the IVT, $g(c) = 0$ for some $c$:
$f(c) = c$.

\textbf{12.} Idealizing the country as the segment
$\intcc{0}{1}$: $f$ sends each place to the place shown
directly beneath it on the crumpled map; crumpling keeps the
map over the country, so $f$ lands in $\intcc{0}{1}$, and it
is \omterm{def:g12:limcont:continuity}{continuous} (paper does not tear). The \omterm{pb:g10:functions:1}{fixed point} $c$ is
the spot lying precisely above its own \omterm{def:g10:functions:function}{image} --- Brouwer's
theorem handles the honest two-dimensional map.

\textbf{13.} $g(0) + g(\pi) = \left(T(0) - T(\pi)\right) +
\left(T(\pi) - T(2\pi)\right) = T(0) - T(2\pi) = 0$. So
$g(0)$ and $g(\pi)$ are opposite (or both zero): either way
$g$, \omterm{def:g12:limcont:continuity}{continuous}, vanishes at some $t^* \in \intcc{0}{\pi}$:
$T(t^*) = T(t^* + \pi)$ --- somewhere on the equator, two
antipodal points agree on the temperature, at every single
moment.

\textbf{14.} Rotating the table by a quarter turn exchanges
the roles of the two diagonals: the pair of legs that pressed
the ground and the pair that (one of them) hovered swap
situations, so the hovering height --- counted with its sign
as ``gap of the distinguished diagonal'' --- takes opposite
signs at $\theta$ and $\theta + \frac\pi2$. \omterm{def:g12:limcont:continuity}{Continuous} ground
makes $h$ \omterm{def:g12:limcont:continuity}{continuous}, and the IVT hands over an angle
$\theta^*$ with $h(\theta^*) = 0$: four legs down. (No such
luck with a rectangular table or stepped ground.)

\textbf{15.} Template: build the \emph{difference} of the two
quantities you want to coincide ($f(x) - x$;
$T(t) - T(t + \pi)$; up-position minus down-position for the
hiker; the signed gap for the table); check it takes values of
both signs (often via a symmetry of the ends); invoke the IVT
on a \omterm{def:g12:limcont:continuity}{continuous} \omterm{def:g11:func:function}{function}. Coincidence forced.

\textbf{16.} Open questions: \emph{how many} solutions
(answered by strict monotonicity and the bijection theorem,
question 7) and \emph{where} they are (answered by \omterm{thm:g12:limcont:ivt}{dichotomy}
or Newton, question 9). The IVT is a pure existence oracle.

\textbf{17.} $3 - \frac{7}{x + 2}$: \omterm{def:g12:limcont:asymptote}{vertical asymptote}
$x = -2$, horizontal $y = 3$; on each side of $-2$ the
\omterm{def:g11:func:function}{function} is \omterm{def:g12:seq:monotonic}{increasing} ($-\frac{7}{x+2}$ increases where
defined); the curve is the reference \omterm{def:g11:func:reference}{hyperbola}
$-\frac7x$ shifted left $2$ and up $3$ --- two branches
hugging the \omterm{def:g12:limcont:asymptote}{asymptote} cross.

\textbf{18.} For $x \neq 1$: $f(x) = x + 1$, so
$\lim_{x \to 1} f = 2$: setting $f(1) = 2$ restores
\omterm{def:g12:limcont:continuity}{continuity}. The discontinuity was a puncture, not a jump ---
fill the single missing point and the curve heals:
\emph{removable}.

\textbf{19.} \omterm{def:g12:limcont:continuity}{Continuous} on every \omterm{def:g10:numbers:interval}{interval} between consecutive
\omterm{def:g10:numbers:sets}{integers}; at each \omterm{def:g10:numbers:sets}{integer} $n$ it jumps from $n - 1$ to $n$.
With $a = 0.9$, $b = 1.1$: $\lfloor a \rfloor = 0 < \frac12 <
1 = \lfloor b \rfloor$, yet $\lfloor x \rfloor = \frac12$ has
no solution --- the jump teleports over $\frac12$, which is
precisely what the IVT forbids \omterm{def:g12:limcont:continuity}{continuous} \omterm{def:g11:func:function}{functions} to do.

\textbf{20.} Local face: value $=$ limit at each point ---
the removable puncture (question 18) violated it curably, the
floor's jumps (question 19) incurably. Global face: on an
\omterm{def:g10:numbers:interval}{interval}, all intermediate values are visited --- the \omterm{def:g11:quad:discriminant}{root}
certificates (6--8), the \omterm{pb:g10:functions:1}{fixed point}, the equator, the table
and the hiker all live here. The slogan's official name: the
\emph{intermediate value theorem}.
\end{solution}
